\section{一次式の積分}\label{節:一次式の積分}

前節の表には $\sin x$ や $\e^x$ はあるが,\,$\sin(3x)$ や $\e^{2x-1}$ はない.では,変数 $x$ のところが一次式に置き換わった
\[
 \int f(ax+b)\dd x\qquad(a\ne0)
\]
はどうなるだろうか.\,$f$ の原始関数 $F$ は表から分かっているとして,その $F$ に帰着させることを目標とする.

まずは $\sin(3x)$ で試してみよう.\,$-\cos(3x)$ を候補にして微分すると
\[
 \frac{\dd}{\dd x}\{-\cos(3x)\}=3\sin(3x)
\]
となる.合成関数の微分によって,内側の $3x$ の微分 $3$ が余分に掛かったのだ.したがって候補を $3$ で割ればよい.

一般の場合も同じである.\,$F'=f$ とすると,合成関数の微分から
\[
 \frac{\dd}{\dd x}F(ax+b)=af(ax+b)
\]
となり,余分な $a$ が掛かる.両辺を $a$ で割って
\[
 \frac{\dd}{\dd x}\left\{\frac1aF(ax+b)\right\}=f(ax+b)
\]
を得る.すなわち
\[
 \int f(ax+b)\dd x=\frac1aF(ax+b)+C
\]
である.\,$a$ で割るのは書式上の約束ではなく,微分したときに現れる $a$ を打ち消すためである.

分母が一次式の $\frac1{ax+b}$ では,\,$\log|x|$ から $\frac1a\log|ax+b|+C$ を得る.この場合は $ax+b=0$ となる $x=-\frac ba$ を除き,その左右の各区間で考える.

\begin{bookpoint}
\textbf{POINT}\par
表の原始関数 $F$ に $ax+b$ を代入し,内側の係数 $a$ で割って $\frac1aF(ax+b)$ とする.
\end{bookpoint}

\begin{example}{内側の微分を確かめる}
次の不定積分を求めよ.
\[
 \int\left(4\e^{2x-1}+\frac3{5x+2}-2\cos(3x)\right)\dd x.
\]
\end{example}

\paragraphbf{方針}
3つの項を $\e^x$,\,$\frac1x$,\,$\cos x$ の $x$ を一次式に置き換えたものとみて,内側の係数 $2,5,3$ でそれぞれ割る.

\paragraphbf{解答}
$2\e^{2x-1}+\frac35\log|5x+2|-\frac23\sin(3x)+C$

\paragraphbf{解法}
指数の $2x-1$,分母の $5x+2$,角の $3x$ の $x$ の係数は,それぞれ $2,5,3$ である.これらで割ると
\begin{align*}
 &\int\left(4\e^{2x-1}+\frac3{5x+2}-2\cos(3x)\right)\dd x\\
 &\qquad=4\cdot\frac12\e^{2x-1}+3\cdot\frac15\log|5x+2|-2\cdot\frac13\sin(3x)+C\\
 &\qquad=2\e^{2x-1}+\frac35\log|5x+2|-\frac23\sin(3x)+C
\end{align*}
となる.ただし,\,$5x+2=0$ となる $x=-\frac25$ は除き,\,$x<-\frac25$ と $x>-\frac25$ の各区間で考える.対数の項を微分すると
\[
 \frac35\cdot\frac5{5x+2}=\frac3{5x+2}
\]
となり,分母の係数 $5$ が打ち消されて元の項に戻る.

\paragraphbf{内側が一次式でない場合}
一次式の場合は内側の微分が定数 $a$ なので,係数を直すだけで済んだ.内側が $x^2+1$ ならば,微分は $2x$ であり,定数ではない.\,$F(x^2+1)$ を $2x$ で割った候補を微分しても,商の微分から余分な項が出てしまう.次章の置換積分では,この $2x$ が積分される式の中にあるかどうかが手がかりになる.
